\documentclass[aspectratio=169, 10pt]{beamer}
\usepackage{beamer_theme_cienciadados2026}

\title[Aula 04 - Linux e Git]{Ciência dos Dados I: Análise Exploratória}
\subtitle{Aula 04: Ambiente Unix/Linux, Linha de Comando e Controle de Versão}
\author{Prof. Dr. Ney Lemke}
\institute{UNESP -- Instituto de Biociências de Botucatu}
\date{Versão 2026}

\begin{document}

\frame{\titlepage}

\begin{frame}{Sumário da Aula}
  \tableofcontents
\end{frame}

\section{A Filosofia Unix}

\begin{frame}{Princípios do Design Unix}
  Formulada por Doug McIlroy e Ken Thompson nos laboratórios Bell:
  \begin{enumerate}
    \item \textbf{Faça uma coisa bem feita}: Cada utilitário resolve uma tarefa específica com excelência.
    \item \textbf{Texto é a interface universal}: A saída de um programa pode ser a entrada de outro.
    \item \textbf{Combine ferramentas simples}: O operador pipe (\texttt{|}) encadeia fluxos contínuos de dados.
  \end{enumerate}
  \begin{block}{Por que o Cientista de Dados deve dominar a CLI?}
    Permite filtrar, inspecionar e pré-processar arquivos de 100 GB no servidor em segundos, sem necessidade de alocar memória RAM com Pandas.
  \end{block}
\end{frame}

\section{Streams e Utilitários Essenciais}

\begin{frame}[fragile]{Comandos de Inspeção Rápida de Dados}
  \begin{itemize}
    \item \texttt{head -n 10 arquivo.csv}: Exibe o cabeçalho e as primeiras linhas.
    \item \texttt{tail -n 10 arquivo.csv}: Inspeciona o final do arquivo (útil para detectar truncamentos).
    \item \texttt{wc -l arquivo.csv}: Conta o número total de linhas.
    \item \texttt{grep -i "termo" log.txt | wc -l}: Filtra padrões e conta correspondências.
    \item \texttt{cut -d ',' -f 1,3 arquivo.csv}: Isola colunas específicas em arquivos delimitados.
  \end{itemize}
\end{frame}

\section{Controle de Versão com Git e GitHub}

\begin{frame}{Reprodutibilidade e Versionamento}
  \begin{columns}
    \begin{column}{0.5\textwidth}
      \begin{itemize}
        \item \textbf{Git}: Sistema distribuído que rastreia snapshots do código ao longo do tempo.
        \item \textbf{Estados}: \textit{Working Directory} $\rightarrow$ \textit{Staging Area} $\rightarrow$ \textit{Repository (.git)}.
        \item \textbf{Branching}: Permite testar modelos e análises sem corromper a linha principal de produção.
      \end{itemize}
    \end{column}
    \begin{column}{0.48\textwidth}
      \begin{block}{Comandos Fundamentais}
        \texttt{git clone <url>} \\
        \texttt{git status} \\
        \texttt{git add .} \\
        \texttt{git commit -m "Mensagem"} \\
        \texttt{git push origin main}
      \end{block}
    \end{column}
  \end{columns}
\end{frame}

\begin{frame}{Resumo da Aula}
  \begin{itemize}
    \item A linha de comando é a primeira linha de defesa no manuseio de dados brutos.
    \item O Git garante que nenhuma análise científica seja perdida ou substituída acidentalmente.
    \item \textbf{Prática}: Comandos de inspeção e contagem executados diretamente das células do Jupyter.
  \end{itemize}
\end{frame}

\end{document}
